\section{Memory and caches}
\label{sec:memory}

Two results above point to memory: the worst prediction of the cost model was
the run with the largest clause database, and skips are cheapest on fullN. The
second curve in Figure~\ref{fig:memory} (left) explains the latter. Reading
memory in order costs 4--8 cycles at any size, and although \inst{full17} uses
31\,MB, almost twice the L2, it creates its clauses one after another and
mostly visits them in that order. Size alone does not make memory slow; size
combined with a random access order does. Estimated from the data structures
(32 bytes per clause object, 4 bytes per literal in a separate block, 8 bytes
of watch entries), the clause data of \inst{php11\_10} takes about 4\,MB, well
inside L2, so in the normal runs memory is not the bottleneck.

\begin{figure}[ht]
\centering
\plotfigure{figures/fig-memory}
\caption{Left: load latency on this machine, for a random pointer chase and an
in-order walk. Right: extra cycles per clause visit beyond the in-L2 cost
model, as clause padding grows the footprint with the search held fixed. The
dashed line marks the 16\,MB L2.}
\label{fig:memory}
\end{figure}

To isolate memory, I added a setting, \texttt{SAT\_CLAUSE\_PAD}, that reserves
unused space after every clause's literals. The literals are untouched, so the
search is identical (every counter was checked), while the footprint grows up
to \numPadFootprint{} times, from 11 to 226\,MB on \inst{php11\_10}. Comparing
each run with the cost model fitted only on in-L2 runs gives the extra cycles
per visit caused by memory alone. They stay near zero up to an estimated
15\,MB of clause data, rise between 15 and 40\,MB, where the latency curve
leaves L2, and reach about 9 cycles at the largest padding, making each visit
\numPadSlowdown\% more expensive (Figure~\ref{fig:memory}, right). That is far
less than one main-memory access, about 450 cycles, because independent visits
overlap their misses.

\begin{table}[ht]
\centering\small
\begin{tabular}{lrrr}
\toprule
region of \texttt{propagate()} & pad 0 & pad 2048 & change \\
\midrule
% generated by make_data.py - do not edit by hand
replacement search & 17.1 & 19.9 & +2.8 \\
unattributed & 4.0 & 6.1 & +2.1 \\
already satisfied & 3.6 & 5.6 & +2.0 \\
reach the clause & 4.6 & 5.7 & +1.1 \\
watch-list walk & 2.1 & 1.9 & -0.1 \\
move the watch & 0.8 & 0.7 & -0.2 \\
\bottomrule

\end{tabular}
\caption{Cycles per clause visit by code region on \inst{php11\_10}, same
search, unpadded and padded. Only regions that read clause literals grow.}
\label{tab:shift}
\end{table}

A profile of the padded run shows where the cost lands
(Table~\ref{tab:shift}): only in code that reads the padded clause data. The
watch-list walk and the watch move touch only memory the padding did not
change, and they do not slow down, which makes them a control group inside the
same measurement. Padding and a large clause database may leave the cache in
different ways. Padding spreads the same hot data over more memory pages; one
plausible explanation of its cost is that the processor runs out of
address-translation entries (the TLB) rather than out of cache, but these
measurements cannot tell TLB misses from cache misses. With clause deletion
turned off, by contrast, \inst{php11\_10} keeps 166{,}675 learnt clauses and
genuinely reads more data, at about 10 extra cycles per visit and 1.70
instructions per cycle. Without miss counts, four signals agree: the latency
curve, the knee at the L2 size, the falling instructions per cycle, and the
cost appearing only where padded data is read. Together they show that memory
becomes more expensive beyond L2, but not which structure is responsible; a
pointer chase that touches one cache line per page would measure the TLB's
reach and separate the two.
